\section{Sharp precision and sample requirements}\label{sec:sharp-precision}

The sharp precision bound separates the information supplied by response labels from that supplied by all covariate-treatment records. The rate and attaining decision are defined once in \cref{obj:def:sharp-rate,obj:def:sharp-decision}. The following result gives their uniform minimax comparison, explicit sample-size regimes, and extension to randomized decisions.

\begin{definitionv}{def:sharp-rate}[Sharp rate \(r_*\)]
The sharp rate at the public class parameters is
\[
r_*(n,m;d,\alpha,\beta,\gamma)
=
\max\left\{
n^{-\gamma/(2\gamma+d)},
[n(n+m)]^{-\gamma/(2\gamma+d+\gamma d/(\alpha+\beta))}
\right\}
=
r_{\mathrm{up}}(n,m),
\]
where \(r_{\mathrm{up}}(n,m)\) is the maximum of the supplied-propensity and unequal-channel rate orders.
\end{definitionv}

\begin{theoremv}{thm:sharp-annotation-frontier}[Sharp annotation frontier]
Fix the covariate dimension \(d\in\mathbb N\), the propensity and control-mean smoothness orders \(\alpha,\beta\), the effect smoothness order \(\gamma\), the Hölder radius \(L\), and the overlap level \(\varepsilon\), satisfying:
\begin{itemize}
\item \textbf{(Dimension.)} \(d\ge1\).
\item \textbf{(Nuisance smoothness.)} \(0<\alpha\le1\) and \(0<\beta\le1\).
\item \textbf{(Effect smoothness.)} \(1\le\gamma\le3\).
\item \textbf{(Radius.)} \(L>1/2\).
\item \textbf{(Overlap.)} \(0<\varepsilon<1/2\).
\end{itemize}
Let \(\mathcal M\) denote the primitive class at these fixed parameters, and let \(r_*\), \(T_*\), and \(R(n,m)\) be the sharp rate, total data-only decision, and minimax risk of \cref{obj:def:sharp-rate,obj:def:sharp-decision,obj:def:minimax}, respectively. There exist real constants \(0<c\le C<\infty\), depending only on the fixed parameters, such that, simultaneously for every \(n,m\in\mathbb N\) with \(n\ge2\), \(T_*\) is Borel measurable and
\[
c\,r_*
\le R(n,m)
\le \sup_{P\in\mathcal M}\mathcal R_{n,m}(T_*,P)
\le C\,r_*.
\]
Moreover, every Borel real-valued decision \(T\) using the original dataset and the independent uniform randomizer satisfies
\[
c\,r_*\le \sup_{P\in\mathcal M}\mathcal R_{n,m}(T,P).
\]

Define the total treatment-record count \(N\), the combined nuisance smoothness \(S\), the exponent denominator \(\Delta\), the smoothness threshold \(S_{\mathrm{crit}}\), the total-record growth exponent \(q_*\), and the supplied-propensity order \(r_o(n)\) by
\[
N=n+m,\qquad S=\alpha+\beta,\qquad
\Delta=2\gamma+d+\frac{\gamma d}{\alpha+\beta},
\]
\[
S_{\mathrm{crit}}=\frac{\gamma d}{2\gamma+d},
\qquad
q_*=\frac{\gamma d}{(\alpha+\beta)(2\gamma+d)},
\qquad
r_o(n)=n^{-\gamma/(2\gamma+d)}.
\]
For these sample sizes, the sharp rate has the branches
\[
r_*=
\begin{cases}
r_o(n), & S\ge S_{\mathrm{crit}},\\
(nN)^{-\gamma/\Delta}, & S<S_{\mathrm{crit}},\quad N\le n^{q_*},\\
r_o(n), & S<S_{\mathrm{crit}},\quad N\ge n^{q_*}.
\end{cases}
\]
Whenever \(N=n^{q_*}\), the two rate formulas coincide:
\[
(nN)^{-\gamma/\Delta}=r_o(n).
\]
\end{theoremv}

The two terms express complementary precision requirements. The supplied-propensity term reflects the labeled-sample requirement for estimating the smooth contrast at a point; the product term reflects the joint requirement on labeled and treatment-record counts when the nuisance functions are unknown. If \(S\ge S_{\mathrm{crit}}\), the labeled count already supports the supplied-propensity order. Equality \(S=S_{\mathrm{crit}}\) belongs to this regime: then \(q_*=1\), and the necessary inequality \(N\ge n\) places every admissible sample pair at or beyond the saturation threshold. If \(S<S_{\mathrm{crit}}\), the exponent \(q_*>1\) creates an intermediate regime \(n\le N\le n^{q_*}\), in which precision depends on the product \(nN\). At \(N=n^{q_*}\), the two terms coincide; for larger total counts, the supplied-propensity term governs the rate.

The lower bound concerns the original records and includes every Borel decision using the independent uniform randomizer. Thus the comparison applies to both deterministic and randomized estimation procedures. The original-record construction in \cref{obj:thm:marked-component-certificate} underlies the sparse-range nuisance converse in \cref{obj:lem:nuisance-frontier-converse}. That converse supplies the product-of-counts lower bound in the range where it governs precision. Combined with the supplied-propensity lower bound below, it yields the sharp comparison across all sample imbalances. The matching upper bound is attained by the publicly tuned decision through \cref{obj:thm:rectangular-upper}.

To interpret saturation, the next result evaluates the same primitive class when the designated propensity function is supplied to the decision. This benchmark identifies the labeled-sample precision against which the value of outcome-free records is measured. Its attaining known-design decision uses
\[
h=n^{-1/(2\gamma+d)},\qquad
T^e(\mathcal D_{n,m},U;e_P)
=\mathrm{clip}_{[-1,1]}\!\left[
\frac1n\sum_{i=1}^n
w_h(X_i)\,r_0^\top r(X_i)
\left\{
\frac{A_iY_i}{e_P(X_i)}-
\frac{(1-A_i)Y_i}{1-e_P(X_i)}
\right\}
\right].
\]
Here \(r_0=r(x_0)\), and \(w_h\) and the degree-two basis \(r\) use the centered localization measure from \cref{obj:synth_3,obj:synth_1}. These formulas are defined for every bandwidth \(h>0\); the supplied-propensity smoother below uses only \(h_n=n^{-1/(2\gamma+d)}\), which lies in \((0,1)\) for \(n\ge2\) and can exceed \(1/2\) for small \(n\). This is a known-design polynomial evaluation average, not an empirical local least-squares fit.

\begin{theoremv}{thm:supplied-propensity}[Supplied-propensity point estimation rate]
Fix public parameters satisfying
\begin{itemize}
\item \textbf{(Dimension.)} \(d\in\mathbb N\) and \(d\ge1\).
\item \textbf{(Nuisance smoothness.)} \(0<\alpha\le1\) and \(0<\beta\le1\).
\item \textbf{(Effect smoothness.)} \(1\le\gamma\le3\).
\item \textbf{(Radius.)} \(L>1/2\).
\item \textbf{(Overlap.)} \(0<\varepsilon<1/2\).
\end{itemize}
Let \(\mathcal M\) be the primitive class at these parameters in \cref{obj:def:primitive-class}, and let \(R^e(n,m)\) be the supplied-propensity minimax absolute-error risk in \cref{obj:def:oracle-risk}. Define the supplied-propensity rate and the evaluation profile by
\[
r_o(n)=n^{-\gamma/(2\gamma+d)},
\qquad
x_0=(1/2,\ldots,1/2).
\]
There exist real constants \(0<c\le C<\infty\), depending only on the fixed public parameters, such that, for every \(n,m\in\mathbb N\) with \(n\ge2\),
\[
c\,r_o(n)\le R^e(n,m)\le C\,r_o(n).
\]

The upper bound is attained by the explicit degree-two smoother under the known uniform design, applied to the inverse-propensity pseudo-outcomes
\[
\frac{A_iY_i}{e_P(X_i)}
-
\frac{(1-A_i)Y_i}{1-e_P(X_i)},
\qquad i=1,\ldots,n,
\]
at bandwidth
\[
h=n^{-1/(2\gamma+d)},
\]
with final clipping to \([-1,1]\). Here the smoother uses its localization weight and coarse polynomial basis at this bandwidth for every \(h>0\), including \(h>1/2\). For every \(P\in\mathcal M\), the resulting decision \(T^e\) is measurable as a function of the dataset and randomizer when supplied with the designated admissible propensity \(e_P\), satisfies
\[
\left|T^e(\mathcal D_{n,m},U;e_P)\right|\le1
\]
for every dataset and randomizer value, and obeys
\[
\mathbb E_{\mathsf E_{n,m}(P)}
\left[
\left|T^e(\mathcal D_{n,m},U;e_P)-\tau_P(x_0)\right|
\right]
\le C\,r_o(n).
\]
The sampling law \(\mathsf E_{n,m}(P)\) is the product law of \(n\) independent labeled records, \(m\) independent outcome-free covariate-treatment records, and an independent uniform randomizer \(U\), as specified in \cref{obj:ass:sample-law}.

The same lower bound holds for the supplied-propensity minimax risk over the same-class subfamily with constant propensity and control mean:
\[
c\,r_o(n)\le
\inf_{T^e}
\sup_{\substack{P\in\mathcal M\\ e_P=\mu_{0,P}=1/2}}
\mathbb E_{\mathsf E_{n,m}(P)}
\left[
\left|T^e(\mathcal D_{n,m},U;e_P)-\tau_P(x_0)\right|
\right],
\]
where the constant-witness equalities hold on the unit cube and the infimum ranges over the supplied-propensity decisions of \cref{obj:def:oracle-risk}.
\end{theoremv}

The upper-bound localization formula uses the same centered cube and uniform localization measure as \cref{obj:synth_3}, evaluated at the single bandwidth \(h_n=n^{-1/(2\gamma+d)}\in(0,1)\), which may exceed \(1/2\). Because the cube has side length \(h_n<1\) and is centered at \(x_0=(1/2,\ldots,1/2)\), it lies inside \([0,1]^d\). The rectangular estimator in \cref{obj:thm:rectangular-upper} retains its stated restriction \(0<h\le1/2\).

With the propensity supplied, inverse-propensity weighting gives a response whose conditional mean is the causal contrast. Local polynomial estimation can therefore target the contrast's smoothness directly, yielding the order \(r_o(n)\). The same-class lower bound with constant propensity and control mean isolates the remaining response information requirement: even within this subfamily, the minimax risk is bounded below by a constant times \(r_o(n)\) for every auxiliary count. Since a supplied-propensity decision can also implement any decision based on the original data, this lower bound provides the labeled-sample floor in \cref{obj:thm:sharp-annotation-frontier}.

The sample requirements for attaining this benchmark can be expressed as a region of admissible counts. For a tolerance \leanref{sym:zeta}{\ensuremath{\zeta}}, the annotation region \leanref{sym:Bstar}{\ensuremath{\mathcal B_*(\zeta)}} records the pairs whose sharp rate lies within that tolerance of the supplied-propensity order.

\begin{definitionv}{def:annotation-region}[Annotation region \(\mathcal B_*(\zeta)\)]
The annotation region at tolerance \(\zeta\) is
\[
\mathcal B_*(\zeta)=
\left\{
(n,m)\in\mathbb N_{\ge2}\times\mathbb N_0:
r_*(n,m;d,\alpha,\beta,\gamma)\le\zeta r_o(n)
\right\},
\]
where \(r_*\) is the sharp rate and \(r_o(n)\) is the supplied-propensity point-estimation order.
\end{definitionv}

For \(\zeta\ge1\), the explicit rate gives the equivalent count condition
\[
(n,m)\in\mathcal B_*(\zeta)
\quad\Longleftrightarrow\quad
n+m\ge \zeta^{-\Delta/\gamma}n^{q_*},
\qquad n\ge2,\quad m\ge0.
\]
In particular, \(\zeta=1\) gives the exact rate-level saturation region. Allowing a fixed finite tolerance gives attainment of the supplied-propensity order up to constants. The complete sequence equivalence, strict-improvement criterion, and target-error conditions are collected in \cref{obj:prop:annotation-threshold} and proved in \cref{sec:acquisition-consequences}.

Attainment of the supplied-propensity order and strict improvement over the labeled-only order impose distinct count requirements. Under \(S<S_{\mathrm{crit}}\), strict improvement occurs exactly when \(m_n/n\to+\infty\), whereas supplied-propensity attainment requires the total count to remain at least a positive constant times \(n^{q_*}\). Consequently, diverging auxiliary-to-labeled ratios can improve the minimax order while leaving precision in the intermediate product-of-counts regime. A bounded auxiliary-to-labeled ratio preserves the labeled-only order up to constants. At and above the smoothness threshold, every auxiliary-size sequence attains the supplied-propensity order.

For acquisition planning, \cref{obj:prop:annotation-threshold} gives two simultaneous requirements. For a target absolute-error scale \(0<\eta<1\), the labeled count must support point-estimation precision at scale \(\eta^{-(2+d/\gamma)}\), and the product of labeled and total counts must support nuisance adjustment at scale \(\eta^{-(2+d/\gamma+d/S)}\). These requirements are necessary up to a constant depending on \((d,\alpha,\beta,\gamma,L,\varepsilon)\) and sufficient up to another such constant for uniform expected absolute error at most \(\eta\). Their rate-level form is exact. Thus increasing either count can satisfy the product requirement, while the labeled-count requirement specifies the response information needed at the target precision.

A one-dimensional example makes the threshold concrete. Set \(d=1\), \(\gamma=1\), and \(\alpha=\beta=1/8\), with any fixed \(L>1/2\) and \(\varepsilon\in(0,1/2)\). Then
\[
S=\frac14,\qquad
S_{\mathrm{crit}}=\frac13,\qquad
\Delta=7,\qquad
q_*=\frac43,
\]
and
\[
r_*(n,m)=\max\left\{n^{-1/3},[n(n+m)]^{-1/7}\right\}.
\]
The labeled-only rate is \(n^{-2/7}\), and the exact branch boundary is \(N=n^{4/3}\). Total counts of order \(n^{4/3}\) therefore attain the supplied-propensity order up to constants. For example, with \(n=1000\), the branch boundary is \(N=10000\), corresponding to \(m=9000\) outcome-free records. These counts describe equality of the rate formulas; the minimax risk retains constants that may depend on \(L\) and \(\varepsilon\). In the same example, the exact rate-level requirements for target scale \(\eta\) are
\[
n\ge\eta^{-3},
\qquad
nN\ge\eta^{-7}.
\]
Taking the labeled count at order \(\eta^{-3}\) consequently requires a total count at order \(\eta^{-4}\), illustrating how response labels and routine covariate-treatment records jointly determine attainable precision.
